\subsection{Nonlinear Stability Theory}
% Lyapunov functions, invariant manifolds, linearization, Hartman-Grobman flavor.

Exact formulas are often unnecessary for proving that scalar trajectories remain ordered or trapped. Differential inequalities provide comparison curves for this purpose.

\begin{definition} \label{def:scalar-upper-lower-solutions} \hlt{Upper and Lower Solutions}\\
Let $U\subseteq\R^2$ be open, $f:U\to\R$, $J\subseteq\R$ be open interval. For $u\in C^1(J,\R)$ with $\Gamma_u\defeq\{(t,u(t)):t\in J\}\subseteq U$:
\begin{enumerate}[label=(\roman*)]
\setlength{\itemsep}{0pt}
\item \hlt{Upper solution}: $u'(t)\geq f(t,u(t))$ $\forall t\in J$; it is \hlt{strict} iff the inequality is strict on $J$.
\item \hlt{Lower solution}: $u'(t)\leq f(t,u(t))$ $\forall t\in J$; it is \hlt{strict} iff the inequality is strict on $J$.
\end{enumerate}
\hlt{Key Insight}: Differential inequalities encode one-sided barriers without solving the ODE.
\end{definition}

\begin{definition} \label{def:scalar-isocline} \hlt{Scalar Isocline}\\
Let $U\subseteq\R^2$ be open, $f:U\to\R$, and $c\in\R$.\\
For $x'=f(t,x)$, the level set $\mathcal{I}_c\defeq\{(t,x)\in U:f(t,x)=c\}$ is the \hlt{isocline of slope $c$}.\\
If $J\subseteq\R$ is open and a portion of $\mathcal{I}_c$ is the graph of $w\in C^1(J,\R)$, then $w'>c$ on $J$ makes $w$ a strict upper solution, while $w'<c$ on $J$ makes $w$ a strict lower solution.\\
\hlt{Key Insight}: An isocline fixes the ODE slope, so only the graph slope remains to be compared.
\end{definition}

\begin{lemma} \label{lem:strict-scalar-comparison} \hlt{Strict Scalar Comparison}\\
Let $U\subseteq\R^2$ be open, $f\in C(U,\R)$, $J\subseteq\R$ open, and $[t_0,T]\subseteq J$ with $t_0<T$.\\
Let $I_+\defeq(t_0,T]$ and suppose $x\in C^1(J,\R)$, $\Gamma_x\subseteq U$, $x'=f(\,\cdot\,,x)$.
\begin{enumerate}[label=(\roman*)]
\setlength{\itemsep}{0pt}
\item If $u\in C^1(J,\R)$ is a strict upper solution and $x(t_0)\leq u(t_0)$, then $x<u$ on $I_+$.
\item If $v\in C^1(J,\R)$ is a strict lower solution and $v(t_0)\leq x(t_0)$, then $v<x$ on $I_+$.
\end{enumerate}
\begin{proof}
\begin{enumerate}[label=(\roman*)]
\setlength{\itemsep}{0pt}
\item Set $\Delta\defeq u-x$. If $\Delta(t)=0$, then $x(t)=u(t)$ and $\Delta'(t)=u'(t)-x'(t)=u'(t)-f(t,u(t))>0$.\\
From $\Delta(t_0)\geq0$, continuity if $\Delta(t_0)>0$, and $\Delta'(t_0)>0$ if $\Delta(t_0)=0$, get $\Delta>0$ on $(t_0,t_0+\varepsilon]$.\\
If $t_1\defeq\inf\{t>t_0:\Delta(t)=0\}\leq T$, then $\Delta>0$ on $(t_0,t_1)$ and $\Delta(t_1)=0$, so $\Delta'(t_1)\leq0$.\\
This contradicts $\Delta'(t_1)>0$. Thus $\Delta>0$ on $I_+$.
\item Set $\Delta\defeq x-v$. If $\Delta(t)=0$, then $x(t)=v(t)$ and $\Delta'(t)=x'(t)-v'(t)=f(t,v(t))-v'(t)>0$.\\
As above, $\Delta(t_0)\geq0$ gives $\Delta>0$ near $t_0^+$.\\
A first zero $t_1\in I_+$ would force $\Delta'(t_1)\leq0$, contradicting $\Delta'(t_1)>0$. Hence $\Delta>0$ on $I_+$.
\end{enumerate}
\hlt{Structural Commentary}: At contact, $\Delta'=+$ forbids crossing.\\
\hlt{Key Insight}: Strict differential inequalities forbid barrier crossing.
\end{proof}
\end{lemma}

\begin{definition} \label{def:locally-lipschitz-uniform-in-time} \hlt{Local State Lipschitz Condition}\\
Let $U\subseteq\R^2$ be open, $f:U\to\R$, and $\mathcal{K}(U)\defeq\{K\subseteq U:K\text{ is compact}\}$.\\
$f$ is \hlt{locally Lipschitz in the state, uniformly in time}\\
iff $\forall K\in\mathcal{K}(U)$, $\exists L_K<+\infty$ such that $\abs{f(t,x)-f(t,y)}\leq L_K\abs{x-y}$ whenever $(t,x),(t,y)\in K$.
\end{definition}

\begin{theorem} \label{thm:scalar-lipschitz-comparison} \hlt{Scalar Lipschitz Comparison}\\
Let $U\subseteq\R^2$ be open, $f\in C(U,\R)$, and $\operatorname{Lip}_x(K,L)$ mean $\abs{f(t,p)-f(t,q)}\leq L\abs{p-q}$ whenever $(t,p),(t,q)\in K$.\\
Assume every compact $K\subseteq U$ has some $L_K<+\infty$ with $\operatorname{Lip}_x(K,L_K)$.\\
Let $I\defeq[t_0,T]\subseteq J$, $J\subseteq\R$ open, $t_0<T$. For $z:J\to\R$, set $\Gamma_z(A)\defeq\{(t,z(t)):t\in A\}$, $R_z\defeq z'-f(\,\cdot\,,z)$.\\
Suppose $x,y\in C^1(J,\R)$, $\Gamma_x(J)\cup\Gamma_y(J)\subseteq U$, and $R_x\leq R_y$ on $I$. Then:
\begin{enumerate}[label=(\roman*)]
\setlength{\itemsep}{0pt}
\item If $x(t_0)\leq y(t_0)$, then $x\leq y$ on $I$.
\item If $t_1\in I$ and $x(t_1)<y(t_1)$, then $x<y$ on $[t_1,T]$.
\end{enumerate}
\begin{proof}
\begin{enumerate}[label=(\roman*)]
\setlength{\itemsep}{0pt}
\item Set $\Delta\defeq x-y$, $S\defeq\{\Delta>0\}\cap I\neq\emptyset$. Since $\Delta(t_0)\leq0$, a component $C$ of $S$ has left endpoint $\alpha$ with $\Delta(\alpha)=0$; choose $t_2\in C$, so $\Delta>0$ on $(\alpha,t_2]$.\\
With $K\defeq\Gamma_x([\alpha,t_2])\cup\Gamma_y([\alpha,t_2])$ and $L=L_K$, $R_x\leq R_y$ gives $\Delta'\leq f(\,\cdot\,,x)-f(\,\cdot\,,y)\leq L\Delta$ on $(\alpha,t_2]$.\\
Thus $(e^{-L(t-\alpha)}\Delta)'\leq0$, so $\Delta(t_2)\leq e^{L(t_2-\alpha)}\Delta(\alpha)=0$, contradicting $t_2\in S$.
\item Set $\delta\defeq y-x$. If $t_2\in(t_1,T]$ is the first zero of $\delta$, then $\delta>0$ on $[t_1,t_2)$.\\
With $K\defeq\Gamma_x([t_1,t_2])\cup\Gamma_y([t_1,t_2])$ and $L=L_K$, by $R_x\leq R_y$, $\delta'\geq f(\,\cdot\,,y)-f(\,\cdot\,,x)\geq-L\delta$ on $[t_1,t_2)$.\\
Hence $(e^{L(t-t_1)}\delta)'\geq0$ and $0=e^{L(t_2-t_1)}\delta(t_2)\geq\delta(t_1)>0$, contradiction.
\end{enumerate}
\hlt{Structural Commentary}: $\Delta'\leq L\Delta$ makes $e^{-Lt}\Delta$ monotone.\\
\hlt{Key Insight}: A nonpositive order defect cannot become positive.
\end{proof}
\end{theorem}

\begin{lemma} \label{lem:autonomous-convergent-state-equilibrium} \hlt{A Convergent Autonomous State Is an Equilibrium}\\
Let $f\in C(\R,\R)$ and let $x\in C^1([t_0,+\infty),\R)$ satisfy $x'=f(x)$. If $x(t)\to x_*\in\R$ as $t\to+\infty$, then:
\begin{enumerate}[label=(\roman*)]
\setlength{\itemsep}{0pt}
\item $x'(t)\to0$.
\item $f(x_*)=0$.
\end{enumerate}
\begin{proof}
By continuity, $x'(t)=f(x(t))\to f(x_*)\defeq L$. If $L>0$, then $x'(t)\geq L/2$ eventually, so $x(t)\to+\infty$ by integration; if $L<0$, then similarly $x(t)\to-\infty$. Both contradict $x(t)\to x_*\in\R$, hence $L=0$.
\begin{enumerate}[label=(\roman*)]
\setlength{\itemsep}{0pt}
\item Since $x'(t)\to L=0$, $x'(t)\to0$.
\item Since $L=f(x_*)$, $f(x_*)=0$.
\end{enumerate}
\hlt{Structural Commentary}: The differential equation converts state convergence into derivative convergence.\\
\hlt{Key Insight}: A nonzero limiting velocity is incompatible with a finite limiting state.
\end{proof}
\end{lemma}

\begin{example} \label{ex:riccati-wedge-comparison} \hlt{A Trapping Wedge for $x'=x^2-t^2$}\\
For $f(t,x)\defeq x^2-t^2$ on $(0,+\infty)\times\R$, set $u(t)\defeq t$ and $v(t)\defeq -t$.\\
Then $u'=1>0=f(t,u)$ and $v'=-1<0=f(t,v)$.\\
If $t_0>0$, $x(t_0)=x_0\in[-t_0,t_0]$, $x'=x^2-t^2$, by Lemma \ref{lem:strict-scalar-comparison}, $-t<x(t)<t$ on forward interval with $t>t_0$.
\end{example}

\begin{proposition} \label{prop:riccati-central-wedge-asymptotics} \hlt{Central-Wedge Convergence for $x'=x^2-t^2$}\\
Let $t_0>0$, $-t_0<x_0<t_0$, and let $x:[t_0,T_+)\to\R$ be the maximal forward solution of $x'=x^2-t^2$, $x(t_0)=x_0$.
\begin{enumerate}[label=(\roman*)]
\setlength{\itemsep}{0pt}
\item $-t<x(t)<t$ for $t\in(t_0,T_+)$.
\item $T_+=+\infty$.
\item $x(t)+t\to0$ as $t\to+\infty$.
\end{enumerate}
\begin{proof}
\begin{enumerate}[label=(\roman*)]
\setlength{\itemsep}{0pt}
\item Let $u(t)\defeq t$ and $v(t)\defeq -t$. Since $u'-(u^2-t^2)=1>0$, $v'-(v^2-t^2)=-1<0$, and $v(t_0)<x_0<u(t_0)$, Lemma \ref{lem:strict-scalar-comparison} implies $v(t)<x(t)<u(t)$ on every compact subinterval of $[t_0,T_+)$.
\item By (i), $x'=(x-t)(x+t)<0$, so $x$ is decreasing.\\
If $T_+<+\infty$, then $-T_+<x(t)\leq x_0$ for $t<T_+$, hence $\ell\defeq\lim\limits_{t\uparrow T_+}x(t)\in[-T_+,x_0]$.\\
Local existence at $(T_+,\ell)$ extends $x$, contradicting maximality; hence $T_+=+\infty$.
\item Fix $T\geq\max\{t_0,2\}$, and set $r(t)\defeq\sqrt{t^2-2}$, $w(t)\defeq-r(t)$ for $t\geq T$. Since $r^2-t^2=w^2-t^2=-2$, $r'=t/r>0>-2$, and $w'=-t/r\geq-\sqrt2>-2$, both $r,w$ are strict upper solutions.\\
Let $d(t)\defeq t-x(t)$. By (i), $d>0$ and $t+x>0$, so $d'=1+t^2-x^2=1+d(t+x)>1$.\\
Since $d(t)\to+\infty$ and $t-r(t)\to0$, choose $T_1\geq T$ with $d(T_1)>T_1-r(T_1)$, i.e. $x(T_1)<r(T_1)$.\\
By strict comparison, $x<r$ for $t>T_1$. If $x(T_1)\leq w(T_1)$, then by strict comparison $x<w$ for $t>T_1$.\\
If $w(T_1)<x(T_1)<r(T_1)$, then while $w<x<r$, $x'<-2$ and $w'\geq-\sqrt2$, hence $(x-w)'\leq-(2-\sqrt2)$.\\
Thus $x$ meets $w$ in finite time and then satisfies $x<w$ thereafter.\\
Therefore eventually $0<x(t)+t<t-r(t)=\frac{2}{t+\sqrt{t^2-2}}\to0$.
\end{enumerate}
\hlt{Structural Commentary}: Two nested isocline barriers squeeze the trajectory against $x=-t$.\\
\hlt{Key Insight}: Every trajectory entering the central wedge is asymptotic to its lower edge.
\end{proof}
\end{proposition}

\begin{proposition} \label{prop:riccati-upper-isocline-blowup} \hlt{Finite-Time Blow-Up Above the Upper Isocline}\\
Let $t_0\geq0$, $x_0\geq\sqrt{t_0^2+2}$, and let $x:[t_0,T_+)\to\R$ be the maximal forward solution of $x'=x^2-t^2$, $x(t_0)=x_0$.
\begin{enumerate}[label=(\roman*)]
\setlength{\itemsep}{0pt}
\item $x(t)>\sqrt{t^2+2}$ and $x'(t)>2$ for $t\in(t_0,T_+)$.
\item $T_+<+\infty$ and $x(t)\to+\infty$ as $t\uparrow T_+$.
\end{enumerate}
\begin{proof}
\begin{enumerate}[label=(\roman*)]
\setlength{\itemsep}{0pt}
\item Set $b(t)\defeq\sqrt{t^2+2}$. Then $b^2-t^2=2$ and $b'=\frac{t}{b}<1<2$, so $b'<b^2-t^2$.\\
Since $x_0\geq b(t_0)$, Lemma \ref{lem:strict-scalar-comparison} implies $b(t)<x(t)$ for $t>t_0$; hence $x'=x^2-t^2>b^2-t^2=2$.
\item Suppose $T_+=+\infty$. By (i), $x(t)>x_0+2(t-t_0)$; choose $t_1\geq t_0$ with $x(t)>\sqrt{2}t>0$ for $t\geq t_1$.\\
Then $x'=x^2-t^2>\frac12x^2$ and $(\frac{1}{x})'=-\frac{x'}{x^2}<-\frac12$ on $[t_1,+\infty)$, so $\frac{1}{x(t)}<\frac{1}{x(t_1)}-\frac12(t-t_1)$.\\
This contradicts $x>0$ for $t>t_1+\frac{2}{x(t_1)}$. Thus $T_+<+\infty$. Also $x'>2$, so $\ell\defeq\lim\limits_{t\uparrow T_+}x(t)\in(x_0,+\infty]$.\\
If $\ell<+\infty$, local existence at $(T_+,\ell)$ extends $x$, contradicting maximality; hence $\ell=+\infty$.
\end{enumerate}
\hlt{Structural Commentary}: Above $b(t)=\sqrt{t^2+2}$, linear escape enters the Riccati regime.\\
\hlt{Key Insight}: Once $x'\geq\frac12x^2$, reciprocal decay forces finite-time escape.
\end{proof}
\end{proposition}

\begin{theorem} \label{thm:riccati-separatrix-classification} \hlt{Separatrix Classification for $x'=x^2-t^2$}\\
For $a\in\R$, let $x_a:[0,T_a)\to\R$ be the maximal forward solution of $x'=x^2-t^2$, $x_a(0)=a$.
\begin{enumerate}[label=(\roman*)]
\setlength{\itemsep}{0pt}
\item $\exists!\,a_*\in[0,\sqrt2]$ such that $T_{a_*}=+\infty$, $t\leq x_{a_*}(t)\leq\sqrt{t^2+2}$ for $t\geq0$, and $x_{a_*}(t)-t\to0$ as $t\to+\infty$.
\item If $a>a_*$, then $T_a<+\infty$ and $x_a(t)\to+\infty$ as $t\uparrow T_a$.
\item If $a<a_*$, then $T_a=+\infty$ and $x_a(t)+t\to0$ as $t\to+\infty$.
\end{enumerate}
\begin{proof} Set $\ell(t)\defeq t$ and $b(t)\defeq\sqrt{t^2+2}$.\\
For $T>0$, let $p_T,q_T$ solve $x'=x^2-t^2$ backward from $p_T(T)=\ell(T)$, $q_T(T)=b(T)$.\\
Under $s=T-t$, the reversed equation is $y'=(T-s)^2-y^2$.\\
Hence $\widehat\ell=T-s$, $\widehat b=\sqrt{(T-s)^2+2}$ satisfy $\widehat\ell'<(T-s)^2-\widehat\ell^2$ and $\widehat b'>(T-s)^2-\widehat b^2$.\\
By Lemma \ref{lem:strict-scalar-comparison}, Theorem \ref{thm:scalar-lipschitz-comparison}, and Theorem \ref{thm:picard-local-existence-uniqueness}, $\ell<p_T<q_T<b$ on $[0,T]$.\\
Set $\alpha_T\defeq p_T(0)$, $\beta_T\defeq q_T(0)$, and $K_T\defeq[\alpha_T,\beta_T]$. Then $a\in K_T$ iff $x_a$ exists on $[0,T]$ with $\ell\leq x_a\leq b$.\
Thus $S\geq T$ implies $K_S\subseteq K_T$, and $K_n\subseteq[0,\sqrt2]$ are nonempty compact nested intervals.
\begin{enumerate}[label=(\roman*)]
\setlength{\itemsep}{0pt}
\item Choose $a_*\in\bigcap\limits_{n\in\N}K_n\subseteq[0,\sqrt2]$. Then $x_{a_*}$ is global, $\ell\leq x_{a_*}\leq b$, and $0\leq x_{a_*}(t)-t\leq b(t)-t=\frac{2}{b(t)+t}\to0$.\\
If $a_1<a_2$ lie in $\bigcap\limits_nK_n$, then $\delta\defeq x_{a_2}-x_{a_1}$ satisfies $\delta(0)>0$ and $\delta'=(x_{a_1}+x_{a_2})\delta\geq0$.\\
Thus $0<a_2-a_1\leq\delta(t)\leq b(t)-t\to0$, impossible, so $\bigcap\limits_nK_n=\{a_*\}$.
\item Let $a>a_*$. Choose $n$ with $a>\beta_n$. By comparison, $x_a>q_n>t$, so $x_a>0$ and $x_a'>0$ while both exist.\\
If $T_a\leq n$, $\ell_a\defeq\lim\limits_{t\uparrow T_a}x_a(t)\in(a,+\infty]$; $\ell_a<+\infty$ contradicts maximality (Theorem \ref{thm:picard-local-existence-uniqueness}), $x_a(t)\to+\infty$.\\
If $T_a>n$, then $x_a(n)>q_n(n)=\sqrt{n^2+2}$; Proposition \ref{prop:riccati-upper-isocline-blowup} implies $T_a<+\infty$ and $x_a(t)\to+\infty$ as $t\uparrow T_a$.
\item Let $a<a_*$. Choose $n$ with $a<\alpha_n$. By comparison, $x_a<p_n<t$.\\
On $[0,T_a\wedge n)$, $x_a<-t$ implies $x_a'>0$, while $x_a\geq-t$ cannot later cross below $-t$; hence $x_a\geq\min\{a,-n\}$.\\
By Theorem \ref{thm:picard-local-existence-uniqueness}, $T_a>n$, $x_a(n)<n$. If $-n<x_a(n)<n$, by Proposition \ref{prop:riccati-central-wedge-asymptotics}, $T_a=+\infty$, $x_a(t)+t\to0$.\\
If $x_a(n)\leq-n$, set $h\defeq x_a+t$; while $h<0$, $h'=x_a^2-t^2+1>1$.\\
Hence boundedness and Theorem \ref{thm:picard-local-existence-uniqueness} imply existence of a first $\tau\geq n$ with $h(\tau)=0$.\\
For $t_2>\tau$ close to $\tau$, strict comparison implies $-t_2<x_a(t_2)<t_2$.\\
Proposition \ref{prop:riccati-central-wedge-asymptotics} implies $T_a=+\infty$ and $x_a(t)+t\to0$.
\end{enumerate}
\hlt{Structural Commentary}: Nested corridors isolate the threshold separating blow-up from lower-edge attraction.\\
\hlt{Key Insight}: The orbit asymptotic to $x=t$ is the separatrix between blow-up and attraction to $x=-t$.
\end{proof}
\end{theorem}
